\chapter{Chuleta}
\label{ch:chuleta}

\section{El laboratorio completo, de un tirón}

Desde \cod{tutorial/manual/examples/lab}, con los binarios compilados. Es la secuencia exacta que se ha
ejecutado para este documento.

\begin{lstlisting}[style=sh]
source env.sh
clock 2026-09-01T09:00:00Z
$BIN/fake-idp serve --addr 127.0.0.1:19400 --issuer $IDP --key idp.key --clock-file clock > idp.log 2>&1 &
$BIN/fake-connector -listen 127.0.0.1:18282 > connector.log 2>&1 &
$BIN/alexandria-server -config server.yaml -clock-file clock > server.log 2>&1 &
sleep 3

# Gobierno
tok --sub ana --admin > admin.token
tok --sub lucia > steward.token
curl -sk -X POST $HUB/v1/namespaces -H "Authorization: Bearer $(cat admin.token)" \
  -d '{"slug":"flexgrid","title":"Flexgrid Energy Data Space"}'
curl -sk -X POST $HUB/v1/namespaces/flexgrid/stewards -H "Authorization: Bearer $(cat admin.token)" \
  -d '{"iss":"http://127.0.0.1:19400","sub":"lucia","label":"Lucia (DSO)"}'

# 1.0.0
clock 2026-09-01T10:00:00Z
publish ../dsl/flexgrid-1.0.0.yaml && approve '{}'

# Cliente
sed "s/HUB_DID_HERE/$(hubdid)/" client.yaml > client.local.yaml
$BIN/alexandria-client -config client.local.yaml -state-dir var/client/state -clock-file clock > client.log 2>&1 &
sleep 3
setrule contract-request; setrule contract-agreement-in; setrule transfer-request
syncat 2026-09-03T12:00:00Z && rows

# 1.1.0, compatible
syncat 2026-09-20T10:00:00Z
publish ../dsl/flexgrid-1.1.0.yaml && approve '{}'
syncat 2026-09-23T12:00:00Z && rows

# 2.0.0, ruptura
syncat 2026-10-01T10:00:00Z
publish ../dsl/flexgrid-2.0.0.yaml && approve '{}'
syncat 2026-10-10T12:00:00Z && rows
post /negotiations/urn:uuid:7d2a9f10-ret-neg-0007/agreement a1-agreement
syncat 2026-11-05T12:00:00Z && rows
post /transfers/request t1-transfer

# Revocación
syncat 2026-11-10T00:00:00Z
curl -sk -X POST "$HUB/v1/namespaces/flexgrid/versions/1.1.0:revoke" \
  -H "Authorization: Bearer $(cat steward.token)" \
  -d '{"justification":"Individual consumer profiling was ruled unlawful by the regulator"}'
sleep 2; clock 2026-11-12T00:00:10Z; sleep 2
post /transfers/request t1-transfer
\end{lstlisting}

\section{\texttt{alexandriac}}

\noindent\begin{tabularx}{\linewidth}{@{}L{9cm}Y@{}}
\toprule
\textbf{Comando} & \textbf{Para qué} \\
\midrule
\texttt{alexandriac validate f.yaml} & Validar, sin escribir nada. \\
\texttt{alexandriac compile -o out f.yaml} & Validar y generar los artefactos. \\
\texttt{… -vendor vendor/\allowbreak{} bundle.yaml} & Resolver \emph{bindings} vendorizados. \\
\texttt{… -previous v1.yaml v2.yaml} & Comprobar el versionado contra la versión anterior. \\
\bottomrule
\end{tabularx}

\section{API del hub}

\begin{longtable}{@{}L{8.8cm}L{6.6cm}@{}}
\toprule
\textbf{Operación} & \textbf{Quién} \\
\midrule
\endhead
\texttt{POST /\allowbreak{}v1/\allowbreak{}namespaces} & Admin \\
\texttt{POST /\allowbreak{}v1/\allowbreak{}namespaces/\allowbreak{}\{ns\}/\allowbreak{}stewards} & Admin \\
\texttt{DELETE /\allowbreak{}v1/\allowbreak{}namespaces/\allowbreak{}\{ns\}/\allowbreak{}stewards/\allowbreak{}\{id\}} & Admin \\
\texttt{GET /\allowbreak{}v1/\allowbreak{}events} & Admin \\
\texttt{POST /\allowbreak{}v1/\allowbreak{}namespaces/\allowbreak{}\{ns\}/\allowbreak{}vocabularies} & Steward \\
\texttt{PUT /\allowbreak{}v1/\allowbreak{}namespaces/\allowbreak{}\{ns\}/\allowbreak{}vocabularies/\allowbreak{}\{nombre\}/\allowbreak{}content} & Steward \\
\texttt{POST /\allowbreak{}v1/\allowbreak{}namespaces/\allowbreak{}\{ns\}/\allowbreak{}vocabularies/\allowbreak{}\{nombre\}:refresh} & Steward \\
\texttt{POST /\allowbreak{}v1/\allowbreak{}namespaces/\allowbreak{}\{ns\}/\allowbreak{}drafts} & Steward \\
\texttt{PUT /\allowbreak{}v1/\allowbreak{}namespaces/\allowbreak{}\{ns\}/\allowbreak{}drafts/\allowbreak{}\{id\}}, con \cod{If-Match} & Steward \\
\texttt{POST /\allowbreak{}v1/\allowbreak{}namespaces/\allowbreak{}\{ns\}/\allowbreak{}drafts/\allowbreak{}\{id\}:validate} & Steward \\
\texttt{POST /\allowbreak{}v1/\allowbreak{}namespaces/\allowbreak{}\{ns\}/\allowbreak{}drafts/\allowbreak{}\{id\}:propose} & Steward \\
\texttt{POST /\allowbreak{}v1/\allowbreak{}namespaces/\allowbreak{}\{ns\}/\allowbreak{}candidates/\allowbreak{}\{id\}:approve} & Steward \\
\texttt{POST /\allowbreak{}v1/\allowbreak{}namespaces/\allowbreak{}\{ns\}/\allowbreak{}candidates/\allowbreak{}\{id\}:reject} & Steward \\
\texttt{POST /\allowbreak{}v1/\allowbreak{}namespaces/\allowbreak{}\{ns\}/\allowbreak{}versions/\allowbreak{}\{v\}:revoke} & Steward \\
\texttt{GET /\allowbreak{}v1/\allowbreak{}namespaces/\allowbreak{}\{ns\}/\allowbreak{}events} & Steward \\
\texttt{GET /\allowbreak{}v1/\allowbreak{}namespaces/\allowbreak{}\{ns\}/\allowbreak{}versions} & Cualquiera: el listado firmado \\
\texttt{GET /\allowbreak{}v1/\allowbreak{}namespaces/\allowbreak{}\{ns\}/\allowbreak{}versions/\allowbreak{}\{v\}} & Cualquiera: el registro firmado \\
\texttt{GET /\allowbreak{}pub/\allowbreak{}…} & Cualquiera \\
\bottomrule
\end{longtable}

\section{API de administración del cliente}

\noindent\begin{tabularx}{\linewidth}{@{}L{5.6cm}Y@{}}
\toprule
\textbf{Ruta} & \textbf{Para qué} \\
\midrule
\texttt{GET /\allowbreak{}status} & Estado, versiones locales con su papel, último listado, DID del hub y revisión de reglas. \\
\texttt{GET /\allowbreak{}versions} & Registros conocidos con su estado y si se conserva el paquete. \\
\texttt{GET /\allowbreak{}rules} & El conjunto de reglas; su revisión es el \cod{ETag}. \\
\texttt{PUT /\allowbreak{}rules}, \texttt{POST /\allowbreak{}rules} & Reemplazar el conjunto o añadir una regla. Con \cod{If-Match}. \\
\cod{GET}, \cod{PUT}, \texttt{DELETE /\allowbreak{}rules/\allowbreak{}\{id\}} & Leer, reemplazar o borrar una regla. \\
\texttt{POST /\allowbreak{}rules:explain} & Qué haría el proxy con una petición, sin efectos. \\
\texttt{GET /\allowbreak{}agreements/\allowbreak{}\{id\}} & Un acuerdo capturado y si es utilizable. \\
\texttt{GET /\allowbreak{}audit/\allowbreak{}reports/\allowbreak{}latest} & Último informe de auditoría pasiva. \\
\texttt{GET /\allowbreak{}healthz}, \texttt{GET /\allowbreak{}readyz} & Vida y disponibilidad. \\
\texttt{GET /\allowbreak{}metrics} & OpenMetrics. \\
\bottomrule
\end{tabularx}

\medskip
La API de administración no tiene autenticación: su \emph{listener} no debe ser alcanzable desde fuera.

\section{Códigos vistos en este tutorial}

\begin{longtable}{@{}L{6.9cm}L{2.5cm}L{6cm}@{}}
\toprule
\textbf{Código} & \textbf{Quién} & \textbf{Cuándo} \\
\midrule
\endhead
\cod{API_UNAUTHENTICATED} & Hub, 401 & Operación protegida sin token. \\
\cod{API_FORBIDDEN_ROLE} & Hub, 403 & Token válido sin el papel necesario. \\
\cod{VOCAB_UNPARSEABLE} & Hub, 422 & Contenido que no es del tipo declarado. \\
\cod{VERSION_BUMP_TOO_LOW} & Compilador & El salto de versión no cubre los cambios. \\
\cod{TERM_RETIRED_WITHOUT_DEPRECATION} & Compilador & Retirar sin haber marcado obsoleto antes. \\
\cod{TERM_REMOVED} & Compilador & Un término ha desaparecido. \\
\cod{PROFILE_VERSION_CHANGED} & Compilador & Misma versión, otro contenido. \\
\cod{GATE_SHACL_ONLY} & Compilador, aviso & Gate obligatorio cubierto solo por \cod{shacl}. \\
\cod{BINDING_UNRESOLVED} & Compilador & Rango sin versión, o falta la copia vendorizada. \\
\cod{DATES_GRACE_ON_COMPATIBLE} & Hub, 422 & Gracia en una versión compatible. \\
\cod{DATES_GRACE_TOO_SHORT} & Hub, 422 & Gracia por debajo de la mínima. \\
\cod{DATES_JUSTIFICATION_REQUIRED} & Hub, 422 & Revocación sin justificación. \\
\cod{DATES_REVOKE_LEAVES_NO_HEAD} & Hub, 422 & Revocar la única versión vigente. \\
\cod{NO_VERSION_IN_FORCE} & Cliente & Aún no hay versión vigente. \\
\cod{VERDICT_VIOLATION} & Cliente & El mensaje no conforma con ninguna versión vigente. \\
\cod{VERDICT_NOT_IN_HEAD} & Cliente & Conforma con una versión anterior aún vigente. \\
\cod{CONTEXT_NOT_IN_PACKAGE} & Cliente & Contexto JSON-LD ajeno al paquete. \\
\cod{AGREEMENT_REVOCATION_SCHEDULED} & Cliente, evento & La versión de un acuerdo va a revocarse. \\
\cod{AGREEMENT_VERSION_REVOKED} & Cliente & Transferencia bajo un acuerdo de versión revocada. \\
\cod{SYNC_UNREACHABLE} & Cliente, evento & El hub no responde. \\
\cod{CLIENT_DEGRADED_REJECT} & Cliente & Listado caducado y gate en \cod{reject}. \\
\cod{CLIENT_FAIL_CLOSED} & Cliente & Listado caducado más de \cod{max_staleness}. \\
\cod{POC_EVAL_UNSUPPORTED_ROLE} & Cliente, aviso & Artefacto \cod{shacl} que no se ejecuta. \\
\bottomrule
\end{longtable}

\section{Lo que conviene recordar}

\begin{enumerate}[leftmargin=*, itemsep=3pt]
  \item \textbf{Un \emph{binding} obligatorio necesita un \cod{json-schema} en cada gate.} Con solo \cod{shacl}
        compila, se publica y no valida nada; el compilador avisa con \cod{GATE_SHACL_ONLY}.
  \item \textbf{Una regla valida un nodo.} Para la política de una petición de contrato, \texttt{select: /\allowbreak{}offer}; y
        entonces esa regla no puede validar la envoltura.
  \item \textbf{El artefacto \cod{shacl} activa la comprobación de contextos.} Si lo enlazas, todo contexto que
        citen los mensajes tiene que estar en el paquete con su \cod{url}, incluido el de DSP.
  \item \textbf{Nada se borra y nada publicado cambia.} Obsoleto, luego retirado; y una errata es una versión
        de parche.
  \item \textbf{Retirar no rompe acuerdos; revocar sí}, con dos días de aviso.
  \item \textbf{Con reloj de fichero, todo espera al reloj}: los tokens, la reemisión del listado y los
        reintentos.
\end{enumerate}

\section{Recompilar este documento}

Las fuentes están en \cod{tutorial/manual/}: \cod{main.tex} y un fichero por capítulo en \cod{chapters/}.

\begin{lstlisting}[style=sh]
cd tutorial/manual
make            # con latexmk, tectonic o pdflatex en el PATH
make docker     # con la imagen texlive/texlive:latest-small
\end{lstlisting}

La imagen \cod{latest-small} no trae \cod{tcolorbox}, \cod{titlesec} ni \cod{xurl}; el objetivo \cod{docker}
los instala con \cod{tlmgr} antes de compilar.
